\documentclass[11pt]{article}
\usepackage[T1]{fontenc}
\usepackage[utf8]{inputenc}
\usepackage{lmodern,amsmath,amssymb,amsthm}
\usepackage[margin=0.9in]{geometry}
\usepackage{microtype}
\usepackage{newunicodechar}
\newunicodechar{∀}{\ensuremath{\forall}}
\newunicodechar{∃}{\ensuremath{\exists}}
\newunicodechar{↔}{\ensuremath{\leftrightarrow}}
\newunicodechar{∧}{\ensuremath{\wedge}}
\newunicodechar{ρ}{\ensuremath{\rho}}
\usepackage[colorlinks=true,linkcolor=blue,citecolor=blue,urlcolor=blue,pdfauthor={P. Graham},pdftitle={Faithful local dynamics on every finite acyclic graph}]{hyperref}
\newtheorem{theorem}{Theorem}
\newtheorem{lemma}[theorem]{Lemma}
\newcommand{\HK}{\operatorname{HK}}
\newcommand{\ev}{\operatorname{ev}}
\newcommand{\End}{\operatorname{End}}
\title{Faithful local dynamics on every finite acyclic graph}
\author{P. Graham\\ \small\url{https://orcid.org/0009-0007-8493-1505}}
\date{November 1, 2026}
\begin{document}
\maketitle
\begin{abstract}
Give each vertex of a finite directed acyclic graph a finite set of states,
and let updating a vertex overwrite its state with a function of the states
at its outgoing neighbors. Collina and D'Andrea showed that such updates
always satisfy the relations of the graph's Hecke--Kiselman monoid, introduced
by Ganyushkin and Mazorchuk, and conjectured that some choice of states and
functions satisfies no other relations. They proved this for complete acyclic
graphs, and Forsberg's effective representations settle further special
classes. We prove the conjecture for every finite directed acyclic
graph: some such system represents the monoid faithfully, so two words act
equally in every system exactly when they are equal in the monoid. Each update
is recorded as a formal expression built from the expressions it reads.
Aragona and D'Andrea's cancellation rules reduce every word to one of their
normal forms. In a normal form, the expressions produced at each vertex grow
strictly and all survive in the final state, so the final state determines
the word up to swaps of adjacent commuting letters. Lean checks the full
theorem.

This work has not yet been reviewed by experts. The proofs, the Lean formalization and the text were produced by AI models (from OpenAI and Anthropic), working under the author's direction.
\end{abstract}
\medskip
\noindent\textbf{MSC 2020:} 20M05, 37B15.\quad\textbf{Keywords:} Hecke--Kiselman monoid, Kiselman's semigroup, local dynamics, acyclic graph.

\section{Introduction}
Kiselman's semigroups~\cite{Kiselman,KM} have a generalization for every
directed graph, the Hecke--Kiselman monoids of Ganyushkin and
Mazorchuk~\cite{GM}. Collina and D'Andrea~\cite{CD} found these monoids in
sequential dynamics on a directed acyclic graph: each vertex holds a state,
and updating a vertex overwrites its state using only the states at its
outgoing neighbors. The updates of every such system satisfy the monoid's
relations, and Collina and D'Andrea conjectured that some finite system
satisfies no others. We prove this for every finite directed acyclic graph.

Let $G=(V,E)$ be a finite directed acyclic graph, allowing $V=\varnothing$.
Write $\HK_G$ for the monoid generated by $V$ with relations
\begin{equation}\label{relations}
 i^2=i,\qquad ij=ji\quad(i\ne j\text{ nonadjacent}),\qquad
 iji=jij=ij\quad(i\to j),
\end{equation}
where $i$ and $j$ are nonadjacent when there is no arrow between them in
either direction. A local system $D$ on $G$ assigns each vertex $i$ a set
$S_i$, always finite and nonempty, and a function
$f_i:\prod_{i\to j} S_j\to S_i$. Its update $U_i$ changes coordinate $i$ to
$f_i((s_j)_{i\to j})$ and leaves the others fixed. In particular, $f_i$ does
not read $s_i$; for a graph with one arrow $i\to j$, updating the sink $j$
assigns a fixed value and updating $i$ applies a function to the state at
$j$, so updating a vertex twice in a row has the same effect as updating it
once. Evaluation $\ev_D:V^*\to\End(\prod_i S_i)$ uses ordinary function
composition: the rightmost letter acts first. Here $V^*$ includes the empty
word, and $[w]$ denotes its image in $\HK_G$.

\begin{theorem}\label{main}
For every finite directed acyclic graph $G$, there is a local system $D$
and an injective monoid homomorphism
$\rho:\HK_G\to\End(\prod_i S_i)$ with $\rho([i])=U_i$ for every vertex.
Moreover, for all words $a,b\in V^*$,
\[
 \ev_A(a)=\ev_A(b)\ \text{for every local system $A$ on $G$}
 \quad\Longleftrightarrow\quad [a]=[b].
\]
\end{theorem}

Collina and D'Andrea call the smallest quotient of $\HK_G$ through which
every evaluation factors the universal dynamics monoid $D(G)$
\cite[Section 8]{CD}; the second clause of Theorem~\ref{main} says
$D(G)\cong\HK_G$. This is their Conjecture 8.1 in the journal version
of~\cite{CD} (Section 8, Conjecture 1 of arXiv:1311.3460v2), first posed in
the 2013 preprint (arXiv:1311.3460v1, Section 9, Conjecture 1). Their
Proposition 2.4 shows that the updates of every local system satisfy
\eqref{relations}, and their Theorem 6.3 implies faithfulness for complete
acyclic graphs. Their Conjecture 8.2, that a particular pointed update
system on the complete graph is an initial object among pointed update
systems, is a different question and is not addressed here. Collina and
D'Andrea report a computer check of Conjecture 8.1 on all graphs with at
most four vertices and on most graphs with five \cite[Section 8]{CD}; the results of~\cite{CD} are part of Collina's
thesis~\cite{Collina}, according to its acknowledgments. In 2022 D'Andrea
reported that a redefined linking operation ``works (experimentally)'' on
all graphs with at most eight vertices \cite[slide 21]{talk}.

The proof uses the elementary cancellations of Aragona and D'Andrea
\cite[Lemma 2.1]{AD}. The words they cannot shorten, their normal forms, are
described in \cite[Remark 2.5]{AD}, and \cite[Theorem 3.1]{AD} shows that two
normal forms of the same element of $\HK_G$ differ by commutations of
nonadjacent letters (see also \cite[Corollary 3.2]{AD}). The argument below shows directly that two normal forms
with the same final state in one particular local system differ in this way,
and does not use \cite[Theorem 3.1]{AD}. It also follows the reconstruction
viewpoint of D'Andrea's lecture~\cite[slides 16, 19]{talk}.

Earlier faithful models cover special classes of graphs. Kiselman's
semigroup is the Hecke--Kiselman monoid of a complete acyclic graph
\cite[Example 12]{GM}, and Kudryavtseva and Mazorchuk proved that it has a
faithful representation by matrices with nonnegative integer entries
\cite[Theorem 34]{KM}; in their construction each generator differs from
the identity matrix in a single column \cite[Section 11.1]{KM}. Forsberg
proved that representations of Hecke--Kiselman monoids ``constructed by
Ganyushkin and Mazorchuk'' are effective, that is, faithful
\cite[abstract]{Forsberg}: for every orientation of a path, the graphs of
type $A_n$ \cite[Corollary 4.6]{Forsberg}; for weighted representations of
the graphs $Z_n$ with arrows $a\to v_i\to b$ \cite[Theorem 5.1]{Forsberg};
and for graphs glued from such pieces at a common source or sink
\cite[Theorem 4.5]{Forsberg}. Passing to the opposite graph
\cite[Proposition 13]{GM}, transposing and reducing modulo a suitable prime
turns these representations into faithful local systems, so Conjecture 8.1
was already settled for these graphs. The Hecke--Kiselman monoids of type
$A_n$ also ``naturally appear as monoids of projection functors''
\cite[Section 1]{Forsberg}: Grensing gives defining relations for monoids of
projection functors on admissible tree quivers \cite[Theorem 3.9]{Gr}, and
Grensing and Mazorchuk develop these realizations further~\cite{GrM1}, showing
for admissible tree quivers that ``the monoid generated by dual projection
functors is the Hecke--Kiselman monoid of the underlying quiver''
\cite[abstract]{GrM2}. None of these results covers arbitrary acyclic graphs.

Recent work counts the elements of the complete-graph monoids
\cite[Theorems 2.1 and 3.2]{DS} and gives faithful representations of the
monoids of directed cycles \cite[abstract]{JLZ}, which fall outside our
hypothesis, and of Catalan monoids \cite[abstract]{LJZ}; none of these is
used here. In sequential dynamical systems
\cite{BR,BMR2,BMR3,MRdm,Reidys,GJL}, \cite[Chapter 7]{MR}, local maps act on
undirected graphs and may read the updated vertex's own state
\cite[Definition 2.1]{GJL}, so the relations~\eqref{relations} need not
hold; these works study the update orders of given systems, and, to our
knowledge, none determines the relations satisfied by all local systems.
Collina and D'Andrea cite four of them~\cite{BR,BMR2,BMR3,MR} and pose
Conjecture 8.1 after them. Uluta\c{s} and K\"oro\u{g}lu~\cite{UK} study the
computational cost of their earlier correspondence~\cite{UK0} between
period-3 orbits of sequential dynamical systems on complete graphs and
ternary codes, which extends Defant's correspondence for period-2 orbits and
binary codes~\cite{Defant}.

\section{Proof overview}
At each vertex we store a formal expression: an initial symbol, or a
constructor applied to the current expressions at its outgoing neighbors
(Section~\ref{sec:terms}). Acyclicity makes the set of expressions at each
vertex finite. Every word equals a normal form, a word in which any two
occurrences of a letter $i$ have both an incoming and an outgoing neighbor
of $i$ between them (Lemma~\ref{reduction}). In a normal form the outgoing
neighbor makes successive expressions at $i$ strictly larger, and the
incoming neighbor copies each overwritten expression at $i$ into a later
expression, where it survives (Lemma~\ref{events}). The final state
therefore determines the expressions produced at each vertex, their order,
and, on each arrow, how the updates at its two ends interleave
(Lemma~\ref{projection}). The projection lemma for partially commutative
words, due to Cori and Perrin and to Duboc~\cite{CP,Duboc}, then recovers
the word up to swaps of adjacent commuting letters (Lemma~\ref{trace}); we
prove it below.

\section{Constructor terms and recovery}\label{sec:terms}
Because the rightmost letter acts first, a word $w$ is executed in the
order of its reverse $w^{\mathrm{rev}}$; we call $w^{\mathrm{rev}}$ the
chronological schedule of $w$. An \emph{event} is an occurrence of a letter
in a chronological schedule.

For each vertex $i$, let $T_i$ be the set of formal expressions consisting
of a symbol $X_i$ and one expression $N_i((t_j)_{i\to j})$ for each tuple of
expressions $t_j\in T_j$ at the outgoing neighbors $j$ of $i$:
\[
 T_i=\{X_i\}\ \sqcup\ \{N_i((t_j)_{i\to j}):t_j\in T_j\}.
\]
Distinct tuples give distinct expressions, and none of them is $X_i$. Since
$G$ is acyclic, these sets are defined recursively from the sinks upward.
At a sink, $T_i$ has exactly two elements, $X_i$ and $N_i$ of the empty
tuple; at other vertices $T_i$ is finite because the finitely many child
sets are. The \emph{constructor system} is the local system with $S_i=T_i$
and $f_i((t_j))=N_i((t_j))$: updating $i$ replaces its expression by $N_i$
of the current expressions at its outgoing neighbors. Start at
$x=(X_i)_{i\in V}$. The term of an event is the expression it produces. The
subterms of an expression are the expression, its entries $t_j$, their
entries, and so on; a subterm $N_j(\dots)$ or $X_j$ belongs to $T_j$, and
we say it is at vertex $j$. Set $|X_i|=0$ and
$|N_i((t_j))|=1+\sum_{i\to j}|t_j|$.

Following Aragona and D'Andrea~\cite{AD}, call a word a \emph{normal form}
if between every two occurrences of any letter $i$ it contains an incoming
neighbor and an outgoing neighbor of $i$; by \cite[Remark 2.5]{AD} these are
exactly the words to which none of their elementary cancellations applies.
It suffices to check consecutive occurrences, and the reverse of a normal
form is a normal form.

\begin{lemma}\label{reduction}
Every word is equal in $\HK_G$ to a normal form.
\end{lemma}
\begin{proof}
Let $m$ be a word and $i$ a vertex. The elementary cancellations
\cite[Lemma 2.1]{AD}, whose convention $aba=ab$ for $a\to b$ agrees
with~\eqref{relations}, state:
\begin{align*}
 imi&=im &&\text{if $m$ contains no incoming neighbor of $i$},\\
 imi&=mi &&\text{if $m$ contains no outgoing neighbor of $i$}.
\end{align*}
Both follow from~\eqref{relations} by induction on the length of $m$. For
the second, write $m=m_1k$. If $k=i$ or $k$ is nonadjacent to $i$, use
$i^2=i$ or $ki=ik$ and the induction hypothesis $im_1i=m_1i$. If $k\to i$,
then $ki=iki$, so $im_1ki=(im_1i)ki=m_1iki=m_1ki$. The first identity is
symmetric. These identities are also proved in the formalization. A
shortest word equal to a given word therefore admits neither shortening,
so it is a normal form.
\end{proof}

\begin{lemma}\label{events}
Let a chronological schedule be a normal form, executed in the constructor
system from $x$. The sizes of successive event terms at a fixed vertex
strictly increase. The noninitial subterms of the final coordinates are
exactly the event terms, and those at vertex $i$ are the event terms at $i$.
\end{lemma}
\begin{proof}
At every stage, each stored entry $t_j$ of a current expression was a
previous value of coordinate $j$. By induction on time: every update weakly
increases the size of its coordinate, and in a normal form strictly
increases it. Indeed, an update at $i$ reads the current expressions at the
outgoing neighbors, while the expression it replaces is $X_i$ or stores
earlier values of the same coordinates, which by induction are no larger.
A first update has positive size. Between consecutive events at $i$, the
normal form gives an outgoing neighbor $j$ that fires; by induction that
event strictly increases the $j$ coordinate, so the next sum defining the
size at $i$ is strictly larger. In a normal form a sink occurs at most once.

Every noninitial subterm was produced at an event: this is an induction
on the execution, since an update merely wraps already present terms.
Conversely consider an event at $i$. If its value is never overwritten,
it survives in the final $i$ coordinate. Otherwise choose the next event
at $i$. Between them an incoming neighbor $k\to i$ fires, storing the
value as an entry. If that later $k$ term is overwritten, repeat the
argument. The chosen event moves strictly forward in the finite schedule,
so eventually a final coordinate contains the original term.
\end{proof}

\begin{lemma}\label{projection}
If two chronological schedules that are normal forms reach the same final
state in the constructor system from $x$, then they have the same letter
counts and the same two-letter projection on every edge.
\end{lemma}
\begin{proof}
By Lemma~\ref{events}, the final state determines the set of event terms
at each vertex. Their strictly increasing sizes order that set, giving
the entire event history at that vertex, including its length. Fix an
edge $i\to j$. Each $i$ event records either $X_j$ or one term in the
ordered $j$ history. These values are all distinct, so the recorded entry
identifies the number of $j$ events preceding that $i$ event. These numbers,
together with the total $j$ count, determine the $\{i,j\}$ projection:
place the required additional $j$ letters before each successive $i$,
then place the remaining $j$ letters at the end.
\end{proof}

\begin{lemma}[projection lemma~\cite{CP,Duboc}]\label{trace}
Two words with equal letter counts and equal two-letter projections for
each adjacent pair of distinct vertices differ by swaps of consecutive
nonadjacent letters.
\end{lemma}
\begin{proof}
Induct on the common length. If the first word starts with $i$, locate
the first $i$ in the second word; it exists by the count equality. Every
letter before it is distinct from and nonadjacent to $i$: adjacency would
contradict equality of the corresponding two-letter projections. Commute
this $i$ to the front and remove the initial $i$ from both words. The
remaining counts and projections still agree, so induction applies.
The zero-length case is immediate.
\end{proof}

\begin{proof}[Proof of Theorem~\ref{main}]
Every local system satisfies~\eqref{relations}. Idempotence follows
because an update does not read its own coordinate. Nonadjacent updates
commute. For $i\to j$, the update at $j$ does not read $i$, and the outer
updates overwrite earlier values; direct substitution gives
$U_iU_jU_i=U_jU_iU_j=U_iU_j$. Thus evaluation induces a homomorphism
from $\HK_G$ in every local system.

Use the constructor system. Suppose two words $a,b$ act equally. Choose
normal forms $a',b'$ equal to them in $\HK_G$ by Lemma~\ref{reduction};
their actions still agree. Evaluate at $x$. The chronological schedules
$(a')^{\mathrm{rev}}$ and $(b')^{\mathrm{rev}}$ are normal forms and reach
the same final state. Lemmas~\ref{projection} and~\ref{trace} relate them
by swaps of adjacent nonadjacent letters. Reversing these swaps relates
$a'$ and $b'$ by the same commutations. Hence $[a]=[b]$, proving
injectivity.

Equality in $\HK_G$ implies equal evaluations in every system by the
factorization already proved. Conversely, equality in every system implies
equality in the constructor system, hence equality in $\HK_G$. For the
empty graph the empty product is a singleton and the same argument applies.
\end{proof}

\section{Formal verification}
Lean checks the theorem \texttt{CD13.Faithful.full\_endpoint}, whose
statement \texttt{CD13.FullEndpoint} is defined in the file
\texttt{CD13/Endpoint.lean} of the accompanying Lean project:
{\small
\begin{verbatim}
def FullEndpoint : Prop :=
  ∀ (V : Type u) [Finite V] (E : V → V → Prop), Acyclic E →
    ∃ (D : LocalSystem.{u, u} E)
      (ρ : HK E →* Function.End (State D)),
      Function.Injective ρ ∧
      (∀ i, ρ (PresentedMonoid.of (HKRel E) i) = update D i) ∧
      (∀ a b : FreeMonoid V,
        (∀ A : LocalSystem.{u, u} E, eval A a = eval A b) ↔
          PresentedMonoid.mk (HKRel E) a = PresentedMonoid.mk (HKRel E) b)
\end{verbatim}
}
\noindent Here $V$ is any finite type, possibly empty, and $E$ is the arrow
relation. The hypothesis \texttt{Acyclic E} says that no vertex is related
to itself by the transitive closure of $E$, that is, $E$ has no directed
cycle. \texttt{HK E} is the monoid presented by the four relations of
\eqref{relations}; a \texttt{LocalSystem} gives each vertex a finite
nonempty coordinate type and a function reading exactly the coordinates at
its outgoing neighbors; \texttt{update} and \texttt{eval} are $U_i$ and
$\ev_A$, with ordinary composition, rightmost letter first. The coordinate
types lie in the same universe as $V$; since every finite set has a copy
there, no generality is lost. There is no hypothesis on the number of
vertices, degrees, transitivity or word length.

The constructor finiteness, cancellation, event-size and subterm arguments,
projection recovery, trace reconstruction and final action are all proved
in Lean. The command
\texttt{\#print axioms CD13.Faithful.full\_endpoint} reports only Lean's
three standard axioms: \texttt{propext} (propositional extensionality),
\texttt{Classical.choice} (choice) and \texttt{Quot.sound} (quotient
soundness); no unproved mathematical
assumption is used. The build also checks the statement's type and
elementary boundary cases. Reproduction uses the accompanying Lean
project's pinned toolchain and Mathlib dependencies; its README gives the
two standard build commands and the resource requirements.

\begin{thebibliography}{99}\small\raggedright
\bibitem{CD} E. Collina and A. D'Andrea,
\emph{A graph-dynamical interpretation of Kiselman's semigroups},
J. Algebraic Combin. 41 (2015), 1115--1132.
\href{https://doi.org/10.1007/s10801-014-0569-7}{doi:10.1007/s10801-014-0569-7}.
Preprint \href{https://arxiv.org/abs/1311.3460v1}{arXiv:1311.3460v1}, 2013,
titled \emph{Sequential dynamical systems on finite acyclic oriented graphs
and Hecke-Kiselman monoids}; revised as
\href{https://arxiv.org/abs/1311.3460v2}{arXiv:1311.3460v2}, May 20, 2014,
whose numbering is that of the journal version, with Conjectures 8.1 and 8.2
numbered Conjectures 1 and 2 in Section 8.
\bibitem{AD} R. Aragona and A. D'Andrea,
\emph{Normal form in Hecke-Kiselman monoids associated with simple oriented graphs},
Algebra Discrete Math. 30 (2020), 161--171, \href{https://doi.org/10.12958/adm1571}{doi:10.12958/adm1571}.
Numbering cited from \href{https://arxiv.org/abs/1904.10745v1}{arXiv:1904.10745v1}, 2019;
in the journal version Lemma 2.1, Remark 2.5 and Theorem 3.1 are Lemma 1,
Remark 3 and Theorem 1.
\bibitem{talk} A. D'Andrea, \emph{An invitation to Hecke-Kiselman monoids}, ANTIPODE 2 lecture,
September 12, 2022, slides 16, 19 and 21.
\href{https://hopfalgb.ulb.be/ANTIPODE2/Presentations/DAndrea.pdf}{Lecture slides}.
\bibitem{Kiselman} C. O. Kiselman,
\emph{A semigroup of operators in convexity theory},
Trans. Amer. Math. Soc. 354 (2002), 2035--2053.
\href{https://doi.org/10.1090/S0002-9947-02-02915-X}{doi:10.1090/S0002-9947-02-02915-X}.
\bibitem{KM} G. Kudryavtseva and V. Mazorchuk,
\emph{On Kiselman's semigroup},
Yokohama Math. J. 55 (2009), 21--46.
\href{https://arxiv.org/abs/math/0511374v1}{arXiv:math/0511374v1}.
\bibitem{GM} O. Ganyushkin and V. Mazorchuk,
\emph{On Kiselman quotients of 0-Hecke monoids},
Int. Electron. J. Algebra 10 (2011), 174--191.
\href{https://arxiv.org/abs/1006.0316v3}{arXiv:1006.0316v3}.
\bibitem{CP} R. Cori and D. Perrin,
\emph{Automates et commutations partielles},
RAIRO Inform. Th\'eor. 19 (1985), 21--32.
\href{https://doi.org/10.1051/ita/1985190100211}{doi:10.1051/ita/1985190100211}.
\bibitem{Duboc} C. Duboc,
\emph{On some equations in free partially commutative monoids},
Theoret. Comput. Sci. 46 (1986), 159--174.
\href{https://doi.org/10.1016/0304-3975(86)90028-9}{doi:10.1016/0304-3975(86)90028-9}.
\bibitem{Forsberg} L. Forsberg,
\emph{Effective representations of Hecke-Kiselman monoids of type $A_n$},
\href{https://arxiv.org/abs/1205.0676v4}{arXiv:1205.0676v4}, February 15, 2017.
\bibitem{Gr} A.-L. Grensing,
\emph{Monoid algebras of projection functors},
J. Algebra 369 (2012), 16--41.
\href{https://doi.org/10.1016/j.jalgebra.2012.07.002}{doi:10.1016/j.jalgebra.2012.07.002},
\href{https://arxiv.org/abs/1203.1943v2}{arXiv:1203.1943v2}.
\bibitem{GrM1} A.-L. Grensing and V. Mazorchuk,
\emph{Categorification of the Catalan monoid},
Semigroup Forum 89 (2014), 155--168.
\href{https://doi.org/10.1007/s00233-013-9510-y}{doi:10.1007/s00233-013-9510-y}.
\bibitem{GrM2} A.-L. Grensing and V. Mazorchuk,
\emph{Finitary 2-categories associated with dual projection functors},
Commun. Contemp. Math. 19 (2017), 1650016.
\href{https://doi.org/10.1142/S0219199716500164}{doi:10.1142/S0219199716500164}.
\href{https://arxiv.org/abs/1501.00095v2}{arXiv:1501.00095v2}.
\bibitem{DS} A. D'Andrea and S. Stella,
\emph{Cardinalities of Kiselman monoids and their parity asymptotics},
\href{https://arxiv.org/abs/2610.14331v1}{arXiv:2610.14331v1}, posted October 9, 2026.
\bibitem{JLZ} Z. F. Jin, Y. F. Luo and W. T. Zhang,
\emph{Combinatorics and faithful representations of cycle monoids},
Semigroup Forum, October 20, 2026.
\href{https://doi.org/10.1007/s00233-026-10687-7}{doi:10.1007/s00233-026-10687-7}.
\bibitem{LJZ} Y. F. Luo, Z. F. Jin and W. T. Zhang,
\emph{Finite Gr\"obner--Shirshov bases and faithful representations for Catalan monoids},
Comm. Algebra 54 (2026), 185--195.
\href{https://doi.org/10.1080/00927872.2025.2522155}{doi:10.1080/00927872.2025.2522155}.
\bibitem{Collina} E. Collina,
\emph{Dynamical systems on graphs and Hecke-Kiselman monoids},
PhD thesis, Sapienza Universit\`a di Roma, 2014.
\bibitem{BR} C. L. Barrett and C. M. Reidys,
\emph{Elements of a theory of computer simulation I: sequential CA over random graphs},
Appl. Math. Comput. 98 (1999), 241--259.
\href{https://doi.org/10.1016/S0096-3003(97)10166-7}{doi:10.1016/S0096-3003(97)10166-7}.
\bibitem{BMR2} C. L. Barrett, H. S. Mortveit and C. M. Reidys,
\emph{Elements of a theory of simulation II: sequential dynamical systems},
Appl. Math. Comput. 107 (2000), 121--136.
\href{https://doi.org/10.1016/S0096-3003(98)10114-5}{doi:10.1016/S0096-3003(98)10114-5}.
\bibitem{BMR3} C. L. Barrett, H. S. Mortveit and C. M. Reidys,
\emph{Elements of a theory of simulation III: equivalence of SDS},
Appl. Math. Comput. 122 (2001), 325--340.
\href{https://doi.org/10.1016/S0096-3003(00)00042-4}{doi:10.1016/S0096-3003(00)00042-4}.
\bibitem{MRdm} H. S. Mortveit and C. M. Reidys,
\emph{Discrete, sequential dynamical systems},
Discrete Math. 226 (2001), 281--295.
\href{https://doi.org/10.1016/S0012-365X(00)00115-1}{doi:10.1016/S0012-365X(00)00115-1}.
\bibitem{MR} H. S. Mortveit and C. M. Reidys,
\emph{An Introduction to Sequential Dynamical Systems},
Universitext, Springer, New York, 2008.
\href{https://doi.org/10.1007/978-0-387-49879-9}{doi:10.1007/978-0-387-49879-9}.
\bibitem{Reidys} C. M. Reidys,
\emph{Sequential dynamical systems over words},
Ann. Comb. 10 (2006), 481--498.
\href{https://doi.org/10.1007/s00026-006-0301-y}{doi:10.1007/s00026-006-0301-y}.
\bibitem{GJL} L. D. Garcia, A. S. Jarrah and R. Laubenbacher,
\emph{Sequential dynamical systems over words},
Appl. Math. Comput. 174 (2006), 500--510.
\href{https://doi.org/10.1016/j.amc.2005.04.101}{doi:10.1016/j.amc.2005.04.101}.
\bibitem{UK} T. Uluta\c{s} and M. E. K\"oro\u{g}lu,
\emph{Counting complexity of period-3 orbits in sequential dynamical systems and their relationship to error-correcting codes over ternary fields},
Symmetry 18 (2026), 1658.
\href{https://doi.org/10.3390/sym18101658}{doi:10.3390/sym18101658}.
\bibitem{UK0} T. Uluta\c{s} and M. E. K\"oro\u{g}lu,
\emph{Period-3 orbits of sequential dynamical systems and their relationship to error-correcting codes over finite fields},
Graphs Combin. 42 (2026), 14.
\href{https://doi.org/10.1007/s00373-025-03008-2}{doi:10.1007/s00373-025-03008-2}.
\bibitem{Defant} C. Defant,
\emph{Binary codes and period-2 orbits of sequential dynamical systems},
Discrete Math. Theor. Comput. Sci. 19(3) (2017), 10.
\href{https://doi.org/10.23638/DMTCS-19-3-10}{doi:10.23638/DMTCS-19-3-10}.
\end{thebibliography}
\end{document}
